%%%%%%%%%%%%%%%%%%%%%%%%%%%%%%%%%%%%%%%%%
%  My documentation report
%  Objetive: Explain what I did and how, so someone can continue with the investigation
%
% Important note:
% Chapter heading images should have a 2:1 width:height ratio,
% e.g. 920px width and 460px height.
%
%%%%%%%%%%%%%%%%%%%%%%%%%%%%%%%%%%%%%%%%%


%----------------------------------------------------------------------------------------
%	PACKAGES AND OTHER DOCUMENT CONFIGURATIONS
%----------------------------------------------------------------------------------------

\documentclass[11pt,fleqn]{book} % Default font size and left-justified equations

\usepackage[top=3cm,bottom=3cm,left=3.2cm,right=3.2cm,headsep=10pt,letterpaper]{geometry} % Page margins

\usepackage{xcolor} % Required for specifying colors by name
\definecolor{ocre}{RGB}{52,177,201} % Define the orange color used for highlighting throughout the book

% Font Settings
\usepackage{avant} % Use the Avantgarde font for headings
%\usepackage{times} % Use the Times font for headings
\usepackage{mathptmx} % Use the Adobe Times Roman as the default text font together with math symbols from the Sym­bol, Chancery and Com­puter Modern fonts
\usepackage{microtype} % Slightly tweak font spacing for aesthetics
\usepackage[utf8]{inputenc} % Required for including letters with accents
\usepackage[T1]{fontenc} % Use 8-bit encoding that has 256 glyphs
\usepackage{amsthm}

% Bibliography
\usepackage{csquotes}
\usepackage[style=alphabetic,sorting=nyt,sortcites=true,autopunct=true,autolang=hyphen,hyperref=true,abbreviate=false,backref=true,backend=biber]{biblatex}
\addbibresource{bibliography.bib} % BibTeX bibliography file
\defbibheading{bibempty}{}

%----------------------------------------------------------------------------------------
%	VARIOUS REQUIRED PACKAGES
%----------------------------------------------------------------------------------------

\usepackage{titlesec} % Allows customization of titles

\usepackage{graphicx} % Required for including pictures
\graphicspath{{Pictures/}} % Specifies the directory where pictures are stored
% \graphicspath{{Plots/}}
\usepackage{lipsum} % Inserts dummy text

\usepackage{tikz} % Required for drawing custom shapes

\usepackage[english]{babel} % English language/hyphenation

\usepackage{enumitem} % Customize lists
\setlist{nolistsep} % Reduce spacing between bullet points and numbered lists

\usepackage{booktabs} % Required for nicer horizontal rules in tables

\usepackage{eso-pic} % Required for specifying an image background in the title page

%----------------------------------------------------------------------------------------
%	MAIN TABLE OF CONTENTS
%----------------------------------------------------------------------------------------

\usepackage{titletoc} % Required for manipulating the table of contents

\contentsmargin{0cm} % Removes the default margin
% Chapter text styling
\titlecontents{chapter}[1.25cm] % Indentation
{\addvspace{15pt}\large\sffamily\bfseries} % Spacing and font options for chapters
{\color{ocre!60}\contentslabel[\Large\thecontentslabel]{1.25cm}\color{ocre}} % Chapter number
{}  
{\color{ocre!60}\normalsize\sffamily\bfseries\;\titlerule*[.5pc]{.}\;\thecontentspage} % Page number
% Section text styling
\titlecontents{section}[1.25cm] % Indentation
{\addvspace{5pt}\sffamily\bfseries} % Spacing and font options for sections
{\contentslabel[\thecontentslabel]{1.25cm}} % Section number
{}
{\sffamily\hfill\color{black}\thecontentspage} % Page number
[]
% Subsection text styling
\titlecontents{subsection}[1.25cm] % Indentation
{\addvspace{1pt}\sffamily\small} % Spacing and font options for subsections
{\contentslabel[\thecontentslabel]{1.25cm}} % Subsection number
{}
{\sffamily\;\titlerule*[.5pc]{.}\;\thecontentspage} % Page number
[] 

%----------------------------------------------------------------------------------------
%	MINI TABLE OF CONTENTS IN CHAPTER HEADS
%----------------------------------------------------------------------------------------

% Section text styling
\titlecontents{lsection}[0em] % Indendating
{\footnotesize\sffamily} % Font settings
{}
{}
{}

% Subsection text styling
\titlecontents{lsubsection}[.5em] % Indentation
{\normalfont\footnotesize\sffamily} % Font settings
{}
{}
{}
 
%----------------------------------------------------------------------------------------
%	PAGE HEADERS
%----------------------------------------------------------------------------------------

\usepackage{fancyhdr} % Required for header and footer configuration

\pagestyle{fancy}
\renewcommand{\chaptermark}[1]{\markboth{\sffamily\normalsize\bfseries\chaptername\ \thechapter.\ #1}{}} % Chapter text font settings
\renewcommand{\sectionmark}[1]{\markright{\sffamily\normalsize\thesection\hspace{5pt}#1}{}} % Section text font settings
\fancyhf{} \fancyhead[LE,RO]{\sffamily\normalsize\thepage} % Font setting for the page number in the header
\fancyhead[LO]{\rightmark} % Print the nearest section name on the left side of odd pages
\fancyhead[RE]{\leftmark} % Print the current chapter name on the right side of even pages
\renewcommand{\headrulewidth}{0.5pt} % Width of the rule under the header
\addtolength{\headheight}{2.5pt} % Increase the spacing around the header slightly
\renewcommand{\footrulewidth}{0pt} % Removes the rule in the footer
\fancypagestyle{plain}{\fancyhead{}\renewcommand{\headrulewidth}{0pt}} % Style for when a plain pagestyle is specified

% Removes the header from odd empty pages at the end of chapters
\makeatletter
\renewcommand{\cleardoublepage}{
\clearpage\ifodd\c@page\else
\hbox{}
\vspace*{\fill}
\thispagestyle{empty}
\newpage
\fi}

%----------------------------------------------------------------------------------------
%	THEOREM STYLES
%----------------------------------------------------------------------------------------

\usepackage{amsmath,amsfonts,amssymb,amsthm} % For math equations, theorems, symbols, etc

\newcommand{\intoo}[2]{\mathopen{]}#1\,;#2\mathclose{[}}
\newcommand{\ud}{\mathop{\mathrm{{}d}}\mathopen{}}
\newcommand{\intff}[2]{\mathopen{[}#1\,;#2\mathclose{]}}
\newtheorem{notation}{Notation}[chapter]

%%%%%%%%%%%%%%%%%%%%%%%%%%%%%%%%%%%%%%%%%%%%%%%%%%%%%%%%%%%%%%%%%%%%%%%%%%%
%%%%%%%%%%%%%%%%%%%% dedicated to boxed/framed environements %%%%%%%%%%%%%%
%%%%%%%%%%%%%%%%%%%%%%%%%%%%%%%%%%%%%%%%%%%%%%%%%%%%%%%%%%%%%%%%%%%%%%%%%%%
\newtheoremstyle{ocrenumbox}% % Theorem style name
{0pt}% Space above
{0pt}% Space below
{\normalfont}% % Body font
{}% Indent amount
{\small\bf\sffamily\color{ocre}}% % Theorem head font
{\;}% Punctuation after theorem head
{0.25em}% Space after theorem head
{\small\sffamily\color{ocre}\thmname{#1}\nobreakspace\thmnumber{\@ifnotempty{#1}{}\@upn{#2}}% Theorem text (e.g. Theorem 2.1)
\thmnote{\nobreakspace\the\thm@notefont\sffamily\bfseries\color{black}---\nobreakspace#3.}} % Optional theorem note
\renewcommand{\qedsymbol}{$\blacksquare$}% Optional qed square

\newtheoremstyle{blacknumex}% Theorem style name
{5pt}% Space above
{5pt}% Space below
{\normalfont}% Body font
{} % Indent amount
{\small\bf\sffamily}% Theorem head font
{\;}% Punctuation after theorem head
{0.25em}% Space after theorem head
{\small\sffamily{\tiny\ensuremath{\blacksquare}}\nobreakspace\thmname{#1}\nobreakspace\thmnumber{\@ifnotempty{#1}{}\@upn{#2}}% Theorem text (e.g. Theorem 2.1)
\thmnote{\nobreakspace\the\thm@notefont\sffamily\bfseries---\nobreakspace#3.}}% Optional theorem note

\newtheoremstyle{blacknumbox} % Theorem style name
{0pt}% Space above
{0pt}% Space below
{\normalfont}% Body font
{}% Indent amount
{\small\bf\sffamily}% Theorem head font
{\;}% Punctuation after theorem head
{0.25em}% Space after theorem head
{\small\sffamily\thmname{#1}\nobreakspace\thmnumber{\@ifnotempty{#1}{}\@upn{#2}}% Theorem text (e.g. Theorem 2.1)
\thmnote{\nobreakspace\the\thm@notefont\sffamily\bfseries---\nobreakspace#3.}}% Optional theorem note

%%%%%%%%%%%%%%%%%%%%%%%%%%%%%%%%%%%%%%%%%%%%%%%%%%%%%%%%%%%%%%%%%%%%%%%%%%%
%%%%%%%%%%%%% dedicated to non-boxed/non-framed environements %%%%%%%%%%%%%
%%%%%%%%%%%%%%%%%%%%%%%%%%%%%%%%%%%%%%%%%%%%%%%%%%%%%%%%%%%%%%%%%%%%%%%%%%%
\newtheoremstyle{ocrenum}% % Theorem style name
{5pt}% Space above
{5pt}% Space below
{\normalfont}% % Body font
{}% Indent amount
{\small\bf\sffamily\color{ocre}}% % Theorem head font
{\;}% Punctuation after theorem head
{0.25em}% Space after theorem head
{\small\sffamily\color{ocre}\thmname{#1}\nobreakspace\thmnumber{\@ifnotempty{#1}{}\@upn{#2}}% Theorem text (e.g. Theorem 2.1)
\thmnote{\nobreakspace\the\thm@notefont\sffamily\bfseries\color{black}---\nobreakspace#3.}} % Optional theorem note
\renewcommand{\qedsymbol}{$\blacksquare$}% Optional qed square
\makeatother

% Defines the theorem text style for each type of theorem to one of the three styles above
\newcounter{dummy} 
\numberwithin{dummy}{section}
\theoremstyle{ocrenumbox}


\newtheorem{theoremeT}[dummy]{Theorem}
\newtheorem{lemma}[dummy]{Lemma}
\newtheorem{observation}[dummy]{Observation}
\newtheorem{proposition}[dummy]{Proposition}
% \newtheorem{definition}[dummy]{Definition}
\newtheorem{claim}[dummy]{Claim}
\newtheorem{fact}[dummy]{Fact}
\newtheorem{assumption}[dummy]{Assumption}

\newtheorem{problem}{Problem}[chapter]
% \newtheorem{exercise}{Exercise}[chapter]
\theoremstyle{blacknumex}
\newtheorem{exampleT}{Example}[chapter]
\theoremstyle{blacknumbox}
\newtheorem{vocabulary}{Vocabulary}[chapter]
\newtheorem{definitionT}{Definition}[section]
\newtheorem{corollaryT}[dummy]{Corollary}
\theoremstyle{ocrenum}

%----------------------------------------------------------------------------------------
%	DEFINITION OF COLORED BOXES
%----------------------------------------------------------------------------------------

\RequirePackage[framemethod=default]{mdframed} % Required for creating the theorem, definition, exercise and corollary boxes

% Theorem box
\newmdenv[skipabove=7pt,
skipbelow=7pt,
backgroundcolor=black!5,
linecolor=ocre,
innerleftmargin=5pt,
innerrightmargin=5pt,
innertopmargin=5pt,
leftmargin=0cm,
rightmargin=0cm,
innerbottommargin=5pt]{tBox}

% Exercise box	  
\newmdenv[skipabove=7pt,
skipbelow=7pt,
rightline=false,
leftline=true,
topline=false,
bottomline=false,
backgroundcolor=ocre!10,
linecolor=ocre,
innerleftmargin=5pt,
innerrightmargin=5pt,
innertopmargin=5pt,
innerbottommargin=5pt,
leftmargin=0cm,
rightmargin=0cm,
linewidth=4pt]{eBox}	

% Definition box
\newmdenv[skipabove=7pt,
skipbelow=7pt,
rightline=false,
leftline=true,
topline=false,
bottomline=false,
linecolor=ocre,
innerleftmargin=5pt,
innerrightmargin=5pt,
innertopmargin=0pt,
leftmargin=0cm,
rightmargin=0cm,
linewidth=4pt,
innerbottommargin=0pt]{dBox}	

% Corollary box
\newmdenv[skipabove=7pt,
skipbelow=7pt,
rightline=false,
leftline=true,
topline=false,
bottomline=false,
linecolor=gray,
backgroundcolor=black!5,
innerleftmargin=5pt,
innerrightmargin=5pt,
innertopmargin=5pt,
leftmargin=0cm,
rightmargin=0cm,
linewidth=4pt,
innerbottommargin=5pt]{cBox}

% Creates an environment for each type of theorem and assigns it a theorem text style from the "Theorem Styles" section above and a colored box from above
\newenvironment{theorem}{\begin{tBox}\begin{theoremeT}}{\end{theoremeT}\end{tBox}}
\newenvironment{exercise}{\begin{eBox}\begin{exerciseT}}{\hfill{\color{ocre}\tiny\ensuremath{\blacksquare}}\end{exerciseT}\end{eBox}}				  
\newenvironment{definition}{\begin{dBox}\begin{definitionT}}{\end{definitionT}\end{dBox}}	
\newenvironment{example}{\begin{exampleT}}{\hfill{\tiny\ensuremath{\blacksquare}}\end{exampleT}}		
\newenvironment{corollary}{\begin{cBox}\begin{corollaryT}}{\end{corollaryT}\end{cBox}}	

%----------------------------------------------------------------------------------------
%	REMARK ENVIRONMENT
%----------------------------------------------------------------------------------------

\newenvironment{remark}{\par\vspace{10pt}\small % Vertical white space above the remark and smaller font size
\begin{list}{}{
\leftmargin=35pt % Indentation on the left
\rightmargin=25pt}\item\ignorespaces % Indentation on the right
\makebox[-2.5pt]{\begin{tikzpicture}[overlay]
\node[draw=ocre!60,line width=1pt,circle,fill=ocre!25,font=\sffamily\bfseries,inner sep=2pt,outer sep=0pt] at (-15pt,0pt){\textcolor{ocre}{R}};\end{tikzpicture}} % Orange R in a circle
\advance\baselineskip -1pt}{\end{list}\vskip5pt} % Tighter line spacing and white space after remark

%----------------------------------------------------------------------------------------
%	SECTION NUMBERING IN THE MARGIN
%----------------------------------------------------------------------------------------

\makeatletter
\renewcommand{\@seccntformat}[1]{\llap{\textcolor{ocre}{\csname the#1\endcsname}\hspace{1em}}}                    
\renewcommand{\section}{\@startsection{section}{1}{\z@}
{-4ex \@plus -1ex \@minus -.4ex}
{1ex \@plus.2ex }
{\normalfont\large\sffamily\bfseries}}
\renewcommand{\subsection}{\@startsection {subsection}{2}{\z@}
{-3ex \@plus -0.1ex \@minus -.4ex}
{0.5ex \@plus.2ex }
{\normalfont\sffamily\bfseries}}
\renewcommand{\subsubsection}{\@startsection {subsubsection}{3}{\z@}
{-2ex \@plus -0.1ex \@minus -.2ex}
{.2ex \@plus.2ex }
{\normalfont\small\sffamily\bfseries}}                        
\renewcommand\paragraph{\@startsection{paragraph}{4}{\z@}
{-2ex \@plus-.2ex \@minus .2ex}
{.1ex}
{\normalfont\small\sffamily\bfseries}}

%----------------------------------------------------------------------------------------
%	HYPERLINKS IN THE DOCUMENTS
%----------------------------------------------------------------------------------------

% For an unclear reason, the package should be loaded now and not later
\usepackage[hyperindex=true, bookmarks=true]{hyperref}
\hypersetup{hidelinks,colorlinks=false,breaklinks=true,urlcolor= ocre,bookmarksopen=false,pdftitle={Title},pdfauthor={Author}}

%----------------------------------------------------------------------------------------
%	CHAPTER HEADINGS
%----------------------------------------------------------------------------------------

% The set-up below should be (sadly) manually adapted to the overall margin page septup controlled by the geometry package loaded in the main.tex document. It is possible to implement below the dimensions used in the goemetry package (top,bottom,left,right)... TO BE DONE

\newcommand{\thechapterimage}{}
\newcommand{\chapterimage}[1]{\renewcommand{\thechapterimage}{#1}}

% Numbered chapters with mini tableofcontents
\def\thechapter{\arabic{chapter}}
\def\@makechapterhead#1{
\thispagestyle{empty}
{\centering \normalfont\sffamily
\ifnum \c@secnumdepth >\m@ne
\if@mainmatter
\startcontents
\begin{tikzpicture}[remember picture,overlay]
\node at (current page.north west)
{\begin{tikzpicture}[remember picture,overlay]
\node[anchor=north west,inner sep=0pt] at (0,0) {\includegraphics[width=\paperwidth]{\thechapterimage}};
%%%%%%%%%%%%%%%%%%%%%%%%%%%%%%%%%%%%%%%%%%%%%%%%%%%%%%%%%%%%%%%%%%%%%%%%%%%%%%%%%%%%%
% Commenting the 3 lines below removes the small contents box in the chapter heading
%\fill[color=ocre!10!white,opacity=.6] (1cm,0) rectangle (8cm,-7cm);
%\node[anchor=north west] at (1.1cm,.35cm) {\parbox[t][8cm][t]{6.5cm}{\huge\bfseries\flushleft \printcontents{l}{1}{\setcounter{tocdepth}{2}}}};
\draw[anchor=west] (5cm,-9cm) node [rounded corners=20pt,fill=ocre!10!white,text opacity=1,draw=ocre,draw opacity=1,line width=1.5pt,fill opacity=.6,inner sep=12pt]{\huge\sffamily\bfseries\textcolor{black}{\thechapter. #1\strut\makebox[22cm]{}}};
%%%%%%%%%%%%%%%%%%%%%%%%%%%%%%%%%%%%%%%%%%%%%%%%%%%%%%%%%%%%%%%%%%%%%%%%%%%%%%%%%%%%%
\end{tikzpicture}};
\end{tikzpicture}}
\par\vspace*{230\p@}
\fi
\fi}

% Unnumbered chapters without mini tableofcontents (could be added though) 
\def\@makeschapterhead#1{
\thispagestyle{empty}
{\centering \normalfont\sffamily
\ifnum \c@secnumdepth >\m@ne
\if@mainmatter
\begin{tikzpicture}[remember picture,overlay]
\node at (current page.north west)
{\begin{tikzpicture}[remember picture,overlay]
\node[anchor=north west,inner sep=0pt] at (0,0) {\includegraphics[width=\paperwidth]{\thechapterimage}};
\draw[anchor=west] (5cm,-9cm) node [rounded corners=20pt,fill=ocre!10!white,fill opacity=.6,inner sep=12pt,text opacity=1,draw=ocre,draw opacity=1,line width=1.5pt]{\huge\sffamily\bfseries\textcolor{black}{#1\strut\makebox[22cm]{}}};
\end{tikzpicture}};
\end{tikzpicture}}
\par\vspace*{230\p@}
\fi
\fi
}
\makeatother % Insert the commands.tex file which contains the majority of the structure behind the template

%----------------------------------------------------------------------------------------
%	Definitions of new commands
%----------------------------------------------------------------------------------------

\def\R{\mathbb{R}}
\newcommand{\cvx}{convex}
\newcommand{\bra}[1]{\langle #1|}
\newcommand{\ket}[1]{| #1\rangle}
\begin{document}

%----------------------------------------------------------------------------------------
%	TITLE PAGE
%----------------------------------------------------------------------------------------

\begingroup
\thispagestyle{empty}
\AddToShipoutPicture*{\put(0,0){\includegraphics[scale=1.25]{esahubble}}} % Image background
\centering
\vspace*{5cm}
\par\normalfont\fontsize{35}{35}\sffamily\selectfont
\textbf{PRINCIPLES OF QUANTUM MECHANICS}\\
{\LARGE R. Shankar}\par % Book title
\vspace*{1cm}
{\Huge Lecture Notes}\par % Author name
\endgroup

%----------------------------------------------------------------------------------------
%	COPYRIGHT PAGE
%----------------------------------------------------------------------------------------

\newpage
~\vfill
\thispagestyle{empty}

%\noindent Copyright \copyright\ 2014 Andrea Hidalgo\\ % Copyright notice

\noindent \textsc{Notes from the book Principles of Quantum Mechanics.}
\vspace{\baselineskip}

\noindent These notes have been taking while self-studying the book Principles of Quantum Mechanics by R. Shankar, second edition.

\noindent \textit{First release, July 2022} % Printing/edition date

%----------------------------------------------------------------------------------------
%	TABLE OF CONTENTS
%----------------------------------------------------------------------------------------

\chapterimage{head1.png} % Table of contents heading image

\pagestyle{empty} % No headers

\tableofcontents % Print the table of contents itself

%\cleardoublepage % Forces the first chapter to start on an odd page so it's on the right

\pagestyle{fancy} % Print headers again

%----------------------------------------------------------------------------------------
%	CHAPTER 1
%----------------------------------------------------------------------------------------

\chapterimage{head2.png} % Chapter heading image
\chapter{Mathematical introduction}
\section{Vector spaces}
\begin{definition}[Group]
    A group is a set $G$ together with a binary operation here denoted "$\cdot$", that combines any two elements  $a$ and $b$ 
    to form an element of $G$, denoted $a\cdot b$, such that the following three requirements, known as group \textit{axioms}, 
    are satisfied
    \begin{enumerate}
        \item Associativity: $\forall\: a,b,c \in G\; (a\cdot b)\cdot c = a\cdot (b\cdot c)$.
        \item Identity: $\exists\: e \in G\;\mid\; \forall a\in G\; e\cdot a=a\cdot e=a$. 
        \item Inverse $\forall a\in G\;\exists\; a^{-1}\in G \mid a^{-1}\cdot a=a\cdot a^{-1}=e$.
    \end{enumerate}
\end{definition}

\begin{remark}
    In the definition of the identity and inverse element, we clearly stated that they commute with any element of the group. This is
    important to state in the defition, because the operation defined on the group is not necessarily commutative.    
\end{remark}

\begin{remark}
    From the assioms it follows that the identity element is unique. If $e$ and $f$ are both identity elements, then $e=e\cdot f=f$.
    
    Similarly, the inverse of an element $a$ of the group is also unique. If $a$ has both $a_1^{-1}$ and $a_2^{-1}$ as inverse, then
    $a_2^{-1} = a_2^{-1}\cdot e = a_2^{-1}\cdot(a\cdot a_1^{-1}) = (a_2^{-1}\cdot a)\cdot a_1^{-1}=e\cdot a_1^{-1}=a_1^{-1}$.
\end{remark}

\begin{definition}[Abelian group]
    A group is \textit{abelian} if the operation defined on it is commutative.
\end{definition}

\begin{definition}[Field]
    A field is a set $F$ together with two binary operations, called addition and multiplication; 
    it is an \textit{abelian} group under addition with $0$ as the additive identity; 
    the \textit{non zero} elements are an \textit{abelian} group under multiplication with $1$ as the multiplicative identity; 
    and multiplication distributes over addition: $\forall a,b\in F\; a\cdot(b+c)=(a\cdot b) + (a\cdot c)$.
\end{definition}

\begin{remark}
    $\R^2$ is a field with the multiplication defined by $(a,b)\cdot(c,d)=(ac-bd, ad+bc)$. This is exactly the multiplication of two
    complex numbers. Indeed, the complex numbers is the way to give $\R^2$ a field structure. For $n>2$ it is not possible to give $\R^n$
    the structure of a field. However, for $n=4$ we can have a pseudo field using the non-commutative quaternion multiplication.
    For $\R^8$ we can give up the associativity of the multiplication and get the non-associative Caley algebra (octanions).
\end{remark}

\begin{definition}[Vector space]
A vector space over a field $F$ is a set $V$ together with two binary operations. 
The first operation, called vector addition assigns to any two vectors $\ket{V}$ and $\ket{W}$ in $V$ a third vector in $V$ which is commonly written as $\ket{v+w}$.
The second operation, called scalar multiplication, assigns to any scalar $a$ in $F$ and any vector $\ket{V}$ in $V$ another vector in $V$, which is denoted $a\ket{V}$.
These two operations must satisfy the eight following axioms:
\begin{enumerate}
    \item Associativity of vector addition:
    $\forall\: \ket{U}, \ket{V}, \ket{W} \in V\; \ket{U}+(\ket{V}+\ket{W}) = (\ket{U}+\ket{V})+\ket{W}$.
    \item Commutativity of vector addition: $\forall\: \ket{U}, \ket{V} \in V\; \ket{U}+\ket{V}=\ket{V}+\ket{U}$.
    \item Identity element of vector addition:
    $\exists\: \ket{0}\in V\;\mid\;\forall \ket{V}\in V\;\ket{V} + \ket{0}= \ket{V}$. 
    \item Inverse elements of vector addition: $\forall \ket{V}\in V\;\exists\; \ket{-V}\in V \mid 
    \ket{V}+\ket{-V}=\ket{0}$.
    \item Compatibility of scalar multiplication with field multiplication:\\
    $\forall\: a,b\in F\: \ket{V}\in V\;a(b\ket{V}) = (ab)\ket{V}$.
    \item Identity element of scalar multiplication: if $1$ is the multiplicative identity in $F$, then $\forall\:\ket{V}\in V\;1\ket{V}=\ket{V}$.
    \item Distributivity of scalar multiplication with respect to vector addition: $\forall\: a\in F\:\ket{U},\ket{V}\in V\: a(\ket{U}+\ket{V})=a\ket{U}+a\ket{V}$.
    \item Distributivity of scalar multiplication with respect to field addition: 
    $\forall\: a,b\in F\:\ket{V}\in V\:(a+b)\ket{V} = a\ket{V}+b\ket{V}$.
\end{enumerate}
Subtraction of two vectors can be defined as $\ket{V}-\ket{W} = \ket{V} + \ket{-w}$.
\end{definition}

\begin{remark}
    From the axioms it follows that 
    \begin{enumerate}
        \item $\ket{0}$ and the inverse are unique. The proof is the same as for the case of the group above.
        \item $0\ket{V}=0$: $\ket{V}=1\ket{V}=(0+1)\ket{V}=0\ket{V}+\ket{V}\Rightarrow 0\ket{V}=0$.
        \item $\ket{-V}=-\ket{V}$
    \end{enumerate}
\end{remark}

\begin{example}
    Consider all functions $f(x)$ defined in an interval $0<x<L$ and such that $f(0)=f(L)=0$. 
    We define scalar multiplication by $a$ simply as $af(x)$ and addition as pointwise addition. These functions form a vector space where 
    $\ket{0}$ is the function which is zero everywhere and the additive inverse of $\ket{f(x)}$ is $\ket{-f(x)}$.
    Notice that if $f(0)=f(L)\neq 0$ then we do not have a vector space.
\end{example}

\subsection{Linear independency}
\begin{definition}[Linear independency]
    A set of vectors $\ket{i}$ is said to be linearly independent if
    $\sum a_i\ket{i}=\ket{0}$
    is satisfied only by the trivial solution $a_i=0$.
\end{definition}

\begin{remark}
    If a set of vectors is linearly dependent, there must exist at least two $a_i\neq 0$. Let's say $a_3\neq 0$, then
    $\ket{3}=-\sum_{i\neq 3}\frac{a_i}{a_3}\ket{i}$. That is, in a set of linearly dependent vectors, at least one of them
    can be written as a linear combination of the other. On the other hand, it is not possible to write any member of the
    linearly independent set in terms of the others.
\end{remark}

\begin{definition}[Dimension of the vector space]
    A vector space has dimension $n$ if it can accommodate a maximum
    of $n$ linearly independent vectors. It will be denoted by $V^n$.
\end{definition}

\begin{example}
    The set of $2 \times 2$ matrices is a four-dimensional vector space. In fact, the following four vectors are linearly indpenedent:
    \begin{equation}
        \begin{array}{cccc}
            \ket{1}=\left[
                \begin{array}{cc}
                    1&0\\
                    0&0
                \end{array}
        \right] &
        \ket{2}=\left[
            \begin{array}{cc}
                0&1\\
                0&0
            \end{array}
        \right] &
        \ket{3}=\left[
            \begin{array}{cc}
                0&0\\
                1&0
            \end{array}
        \right] &
        \ket{4}=\left[
            \begin{array}{cc}
                0&0\\
                0&1
            \end{array}
        \right]
        \end{array}
    \end{equation} 
    Notice that any other matrix can be written as a linear combination of these four matrices. We have seen that in a set of linearly
    independent vectors, no vector can be written as a linear combination of the others. This means that the four matrices given above 
    represent the largest set of linearly independent vectors, and the vector space of $2\times 2$ matrices has dimension $n=4$.
\end{example}

\begin{definition}[Basis]
    A set of $n$ linearly independent vectors in an $n$-dimensional space is called a \textit{basis}.
\end{definition}

\begin{definition}[Component]
    Given a basis $\ket{i}$, any vector $\ket{V}$ can be written as a linear combination of the vectors $\ket{i}$. The coefficients of the
    linear combination are called the \textit{components} of $\ket{V}$ in the basis $\ket{i}$.
\end{definition}

\begin{theorem}
    The expansion of a vector in a given basis is unique.
\end{theorem}
\begin{proof}
    Let's suppose that there exist two different expansions:
    \begin{equation*}
        \ket{V}=\sum a_i\ket{i} = \sum a'_i\ket{i}
    \end{equation*}
    It follows
    \begin{equation*}
        \sum a_i\ket{i} - \sum a'_i\ket{i} =  \sum (a_i-a'_i)\ket{i} = \ket{0}
    \end{equation*}
    Since $\ket{i}$ are linearly independent, this means that $a_i=a'_i$ and the expansion is unique.
\end{proof}
\end{document}